\providecommand{\SysName}{CARE}

\section{Related Work}
\label{sec:related-work}

\noindent\textbf{Software refactoring and refactoring safety.}
Classical refactoring work frames refactoring as behavior-preserving
restructuring, from Fowler's catalog to surveys and design-smell detection
\cite{fowler2018refactoring,mens2004survey,marinescu2004detection}. Empirical
studies show that developers refactor to manage complexity, but that such
changes remain costly and regression-prone in both industrial and open-source
settings \cite{kim2014empirical,vassallo2019large}. Search-based,
measurement-oriented, large-scale, and history-based studies further examine
how refactoring opportunities are selected, evaluated, and detected
\cite{o2008search,o2016measuring,golubev2021one,tsantalis2018accurate}. Recent
learning-based surveys highlight the promise of neural refactoring while still
emphasizing correctness concerns \cite{nyirongo2024survey}. \SysName{} follows
this safety motivation, but focuses on security-sensitive C/C++ refactoring and
uses validation gates rather than relying on the LLM alone; it is especially
close in spirit to work on making refactorings safer with generated tests
\cite{soares2010making}.

\noindent\textbf{Program analysis and security validation.}
\SysName{} draws on classic data-flow, slicing, SSA, control-dependence, and
interprocedural graph-reachability analyses to construct compact program
context \cite{allen1976program,weiser1984program,cytron1991efficiently,
reps1995precise}. Dynamic, concolic, and symbolic execution systems show how
low-level program behavior can be explored beyond purely syntactic reasoning
\cite{godefroid2005dart,sen2005cute,cadar2008exe,cadar2008klee}. Security
validation for C/C++ is also informed by industrial static analysis, dynamic
instrumentation, sanitizers, spatial and temporal memory-safety enforcement,
and control-flow integrity
\cite{bessey2010few,nethercote2007valgrind,serebryany2012addresssanitizer,
stepanov2015memorysanitizer,nagarakatte2009softbound,nagarakatte2010cets,
duck2018effectivesan,abadi2009control}. Unlike full bug finders or complete
verifiers, \SysName{} uses these ideas as fail-closed checks for generated
refactorings, rejecting patches that introduce obvious leaks, double releases,
unchecked allocations, weakened checks, changed return conventions, or altered
API-visible side effects.

\noindent\textbf{Automated program repair.}
Automated program repair is related because it also generates source-level
patches and validates candidates against executable or semantic criteria.
Generate-and-validate systems explore randomized search, repair templates,
learned repair patterns, staged repair, and patch-space analysis
\cite{le2011genprog,qi2014strength,kim2013automatic,long2015staged,
long2016automatic,long2016analysis}, while semantic repair systems use
constraints or symbolic reasoning to synthesize more targeted fixes
\cite{nguyen2013semfix,mechtaev2015directfix,mechtaev2016angelix}. Benchmark
suites further motivate reproducible evaluation and human-review-like patch
assessment \cite{just2014defects4j,madeiral2019bears,saha2018bugs,
karampatsis2020often,ye2021comprehensive}. \SysName{} differs from APR because
it starts from refactoring opportunities rather than failing tests: a candidate
must preserve intended behavior, cleanup order, validation semantics, and
security-relevant side effects, not merely make tests pass.

\noindent\textbf{Neural and LLM-based code transformation.}
Neural code models provide the foundation for LLM-assisted refactoring.
Pre-trained representation and generation models encode code, natural
language, identifiers, and sometimes data-flow information
\cite{feng2020codebert,guo2020graphcodebert,ahmad2021unified,wang2021codet5},
while large generative code models demonstrate nontrivial program synthesis
ability \cite{chen2021evaluating,li2022competition}. Recent empirical work
studies the potential of LLMs for automated software refactoring
\cite{liu2024empirical}. \SysName{} adopts LLMs as patch proposal engines, but
wraps them in program-context construction, conservative planning, iterative
repair, and correctness/security validation. The resulting architecture is a
guarded synthesis loop rather than direct LLM code generation.
